\documentclass[12pt,a4paper]{article}
\usepackage[a4paper,margin=2.2cm]{geometry}
\usepackage{fontspec}
\usepackage{polyglossia}
\setdefaultlanguage{english}
\setmainfont{TeX Gyre Termes}
\usepackage{enumitem}
\usepackage{setspace}
\setlength{\parindent}{0pt}
\setlength{\parskip}{7pt}
\setstretch{1.2}
\begin{document}

\begin{flushright}
\textbf{Prof. Babiker M. Osman}\\
Cairo\\
Al Nahda University\\
WhatsApp: 00249965428114\\
Mobile: 00201031628478\\
Email: babikerosman@yahoo.com
\end{flushright}

\begin{flushleft}
26 August 2026
\end{flushleft}

\vspace{1em}
\begin{center}
\Large\textbf{To Whom It May Concern}
\end{center}

Muhanad was born in the Kingdom of Saudi Arabia in 2006. He developed viral encephalitis at approximately one and a half years of age, during a holiday visit to Sudan.

Unfortunately, his condition was not initially diagnosed promptly and accurately. He was subsequently seen at a paediatric clinic in Sudan, where he was diagnosed with hypotonic cerebral palsy.

Brain magnetic resonance imaging (MRI) showed diffuse cerebral atrophy and enlargement of the cisterna magna, with no other focal changes.

A subsequent MRI showed cerebral atrophy associated with periventricular leukomalacia, together with a degree of cerebellar atrophy.

These features suggest a microangiopathic state, probably resulting from previous hypoxia or vascular injury. This is consistent with the previous illness and the prolonged untreated convulsions that occurred in El Obeid.

In the Kingdom of Saudi Arabia, Muhanad was generally diagnosed with Lennox--Gastaut Syndrome (LGS), a condition characterized by multiple seizure types, developmental delay, impaired cognitive abilities, and resistance to conventional treatments.

The most recent diagnosis, made by Saudi paediatric neurologists, is:

\begin{center}
\textbf{Cerebral palsy, followed by Acute Necrotizing Encephalopathy of Childhood}
\end{center}

At present, Muhanad is unable to sit or roll onto his side, although he can move his arms and legs. The performance of the right side is noticeably better than that of the left.

He receives F-100 formula through a nasogastric feeding tube. His medications include anticonvulsants, vitamins, and prophylactic antibiotics, as follows:

\begin{enumerate}[label=\arabic*. ,leftmargin=*,itemsep=3pt]
\item Paracetamol (400 mg three times daily).
\item Ibuprofen (200 mg three times daily).
\item Septazole (one tablet twice daily).
\item Flagyl (500 mg twice daily).
\item Phenobarbital (50 mg twice daily).
\item Clonazepam (25.0 mg three times daily).
\item Levetiracetam (900 mg twice daily).
\item Lamotrigine (50 mg twice daily).
\item Domperidone (10 mg twice daily).
\item Esomeprazole (20 mg twice daily).
\item Alverine Spasm (alverine + dimethicone) (one tablet twice daily).
\item Enterogermina (probiotic) (one ampoule twice daily).
\item Enemax (once every two days).
\item Coli-Urenal (choline + hexamine + piperazine) (twice daily).
\item F-100 formula milk (350 mL twice daily).
\item Oral rehydration salts (ORS) sachets (one sachet as needed).
\item Multivitamin syrup (Sansovit with iron) (15 mL once daily).
\item Calcium (Osteocare syrup) (10 mL once daily).
\end{enumerate}

\vspace{1em}
With sincere regards,

\textbf{Prof. Babiker M. Osman}\\
Academic and Researcher in Biomedical Sciences

\end{document}
